
\bta{2007A}
\begin{enumerate}
	%\renewcommand{\labelenumi}{\arabic{enumi}.}
	% A(\Alph) a(\alph) I(\Roman) i(\roman) 1(\arabic)
	%设定全局标号series=example	%引用全局变量resume=example
	%[topsep=-0.3em,parsep=-0.3em,itemsep=-0.3em,partopsep=-0.3em]
	%可使用leftmargin调整列表环境左边的空白长度 [leftmargin=0em]
	\item
（$30'$）在一维无限深方势阱$(0<x<a)$中运动的粒子受到微扰
$$H'(x)=\begin{cases}0,&0<x<\frac{a}{3},\frac{2a}{3}<x<a\\
-V_1,&\frac{a}{3}<x<\frac{2a}{3}\end{cases}$$
作用。试求基态能量的一级修正。


\banswer{
	无扰基态为 $\phi_1(x)=\sqrt{2/a}\sin(\pi x/a)$。一级能量修正为
	\begin{align*}
		E_1^{(1)}
		&=\langle\phi_1|H'|\phi_1\rangle
		=-\frac{2V_1}{a}\int_{a/3}^{2a/3}\sin^2\frac{\pi x}{a}\,dx\\
		&=\boxed{-V_1\left(\frac13+\frac{\sqrt3}{2\pi}\right)}.
	\end{align*}
}


\item 
（$30'$）粒子在势场$V(x)$中运动并处于束缚定态$\psi_n(x)$中。试证明粒子所受势场的作用力的平均值为$ 0  $。


\banswer{
	由 Ehrenfest 定理
	\[
		\frac{d\langle p\rangle}{dt}
		=-\left\langle\frac{dV}{dx}\right\rangle=\langle F\rangle.
	\]
	束缚定态的波函数只随时间获得整体相位，因而 $\langle p\rangle$ 不随时间改变，$d\langle p\rangle/dt=0$。所以
	\[
		\boxed{\langle F\rangle
		=-\int_{-\infty}^{\infty}\psi_n^*(x)V'(x)\psi_n(x)\,dx=0}.
	\]
	该结论要求波函数为正常的束缚定态，使分部积分时的边界项消失。
}


\item 
（$30'$）
\begin{enumerate}
	%\renewcommand{\labelenumi}{\arabic{enumi}.}
	% A(\Alph) a(\alph) I(\Roman) i(\roman) 1(\arabic)
	%设定全局标号series=example	%引用全局变量resume=example
	%[topsep=-0.3em,parsep=-0.3em,itemsep=-0.3em,partopsep=-0.3em]
	%可使用leftmargin调整列表环境左边的空白长度 [leftmargin=0em]
	\item
考虑自旋为$\frac{1}{2}$的系统。试在$(\hat{S}^2,\hat{S}_z)$表象中求算符$A\hat{S}_y+B\hat{S}_z$的本征值及归一化的本征态。其中$\hat{S}_y,\hat{S}_z$是自旋角动量算符，而$A,B$为实常数。

\item 
假定此系统处于以上算符的一个本征态上，求测量$\hat{S}_y$得到结果为$\frac{\hbar}{2}$的概率。
\end{enumerate}


\banswer{
	令 $C=\sqrt{A^2+B^2}$。在 $S_z$ 表象中
	\[
		AS_y+BS_z=\frac\hbar2
		\begin{pmatrix}B&-iA\\iA&-B\end{pmatrix}.
	\]
	\begin{enumerate}
		\item 特征方程为
		\[
		 \det\left[\frac\hbar2\begin{pmatrix}B&-iA\\iA&-B\end{pmatrix}
		             -\lambda I\right]
		 =\lambda^2-\frac{\hbar^2}{4}(A^2+B^2)=0,
		\]
		故本征值为
		\[
			\boxed{\lambda_\pm=\pm\frac\hbar2C}.
		\]
		当 $C+B\ne0$ 时，一组归一化本征态可取
		\[
			|+\rangle=\frac1{\sqrt{2C(C+B)}}
			\begin{pmatrix}C+B\\iA\end{pmatrix},
			\qquad
			|-\rangle=\frac1{\sqrt{2C(C+B)}}
			\begin{pmatrix}iA\\C+B\end{pmatrix}.
		\]
		求正本征矢时，第二行给出 $iAu=(C+B)v$；选择
		$(u,v)=(C+B,iA)$，其模平方为 $(C+B)^2+A^2=2C(C+B)$，
		即得到上式。负本征矢由正交性确定。
		若 $C>0$ 而 $C+B=0$，则 $A=0,B<0$，直接选 $S_z$ 的两个本征态。
		若 $C=0$，原算符为零，任意归一化自旋态都是零本征值的本征态。

		\item 该算符等价于沿单位方向
		$\vec n=(0,A/C,B/C)$ 的自旋投影乘以 $C$（$C>0$）。
		$S_y=+\hbar/2$ 的投影算符为 $\Pi_{+y}=(I+\sigma_y)/2$，
		而 $|\pm\rangle$ 中 $\langle\sigma_y\rangle=\pm A/C$，所以概率分别为
		\[
			\boxed{P_y^{(\pm)}=\frac12\left(1\pm\frac{A}{C}\right)}.
		\]
	\end{enumerate}
}


\newpage
\item 
（$30'$）
两个无相互作用的粒子置于一维无限深势阱中$(0<x<a)$中，对下列两种情况写出两粒子体系具有的两个最低总能量，相应的简并度以及上述能级对应的所有二粒子波函数。
\begin{enumerate}
	%\renewcommand{\labelenumi}{\arabic{enumi}.}
	% A(\Alph) a(\alph) I(\Roman) i(\roman) 1(\arabic)
	%设定全局标号series=example	%引用全局变量resume=example
	%[topsep=-0.3em,parsep=-0.3em,itemsep=-0.3em,partopsep=-0.3em]
	%可使用leftmargin调整列表环境左边的空白长度 [leftmargin=0em]
	\item
两个自旋为$\frac{1}{2}$的可区分粒子；

\item 
两个自旋为$\frac{1}{2}$的全同粒子。

\end{enumerate}


\banswer{
	记单粒子空间本征函数为 $\phi_n(x)$，能量为
	$\varepsilon_n=n^2\varepsilon_1$，并以
	$\chi_\uparrow,\chi_\downarrow$ 表示自旋态。
	\begin{enumerate}
		\item 对两个可区分粒子：
		\begin{itemize}
			\item 最低总能量为 $2\varepsilon_1$。空间态为
			$\phi_1(x_1)\phi_1(x_2)$，乘以四个自旋直积态
			$\chi_{s_1}(1)\chi_{s_2}(2)$（$s_1,s_2=\uparrow,\downarrow$），故简并度为 $4$。
			\item 次低总能量为 $\varepsilon_1+\varepsilon_2=5\varepsilon_1$。空间态可为
			$\phi_1(x_1)\phi_2(x_2)$ 或 $\phi_2(x_1)\phi_1(x_2)$，各自又可乘四个自旋直积态，故简并度为 $8$。
		\end{itemize}

		\item 对两个全同自旋 $1/2$ 粒子，完整波函数必须反对称。记自旋单态为 $\chi_{00}$，三个三重态为 $\chi_{1m}$。
		\begin{itemize}
			\item 最低能量仍为 $2\varepsilon_1$，但只允许
			\[
				\phi_1(x_1)\phi_1(x_2)\chi_{00},
			\]
			故简并度为 $1$。
			\item 次低能量为 $5\varepsilon_1$。全部允许态为
			\[
				\frac{\phi_1(x_1)\phi_2(x_2)+\phi_2(x_1)\phi_1(x_2)}{\sqrt2}\chi_{00},
			\]
			以及
			\[
				\frac{\phi_1(x_1)\phi_2(x_2)-\phi_2(x_1)\phi_1(x_2)}{\sqrt2}\chi_{1m},
				\qquad m=1,0,-1.
			\]
			故简并度为 $1+3=4$。
		\end{itemize}
	\end{enumerate}
}



\item 
（$30'$）一个质量为$m$的粒子被限制在$r=a$和$r=b$的两个不可穿透的同心球面之间运动。不存在其它势，求粒子的基态能量和归一化波函数。（1996 年第四题） 

\banswer{
假定 $0<a<b$。令 $\psi=R_l(r)Y_{lm}(\theta,\varphi)$、$u_l=rR_l$，则
\[
 \left[-\frac{\hbar^2}{2m}\frac{d^2}{dr^2}
       +\frac{\hbar^2l(l+1)}{2mr^2}\right]u_l=Eu_l,
 \qquad u_l(a)=u_l(b)=0.
\]
所有 $l$ 的径向定义区间及端点条件相同，而离心势非负，因此最低能量在 $l=0$ 通道取得。
此时 $u''+k^2u=0$。由左端条件写成 $u=C\sin[k(r-a)]$，右端条件给出
$k(b-a)=n\pi$（$n=1,2,\ldots$）。基态取 $n=1$，故
\[
 \boxed{E_0=\frac{\pi^2\hbar^2}{2m(b-a)^2}}.
\]
球谐函数取归一化的 $Y_{00}=1/\sqrt{4\pi}$，三维归一化就等价于
\[
 1=\int_a^b|R_0|^2r^2dr=|C|^2\int_a^b
       \sin^2\frac{\pi(r-a)}{b-a}dr=|C|^2\frac{b-a}{2}.
\]
所以全波函数为
\[
 \boxed{\psi_0(r,\theta,\varphi)=\begin{cases}
 \dfrac{1}{\sqrt{2\pi(b-a)}}\dfrac{\sin[\pi(r-a)/(b-a)]}{r},&a<r<b,\\[6pt]
 0,&\text{其他区域}.
 \end{cases}}
\]
其中 $1/r$ 来自 $R=u/r$，不能把径向函数直接当作一维正弦函数。
}



\end{enumerate}
